The primitive coproduct of an even polynomial generator already determines
which tensor product a finite-height calculation needs.  Give $x$ charge
$\alpha_i$, hence height one, and use the ordinary Hopf algebra $\kk[x]$.
Then
\begin{equation}
 \Delta(x^2)=x^2\otimes1+2x\otimes x+1\otimes x^2.
 \label{polynomial-cross-term}
\end{equation}
The algebra projection $p:\kk[x]\to\kk[x]/(x^2)$ kills $x^2$, but
$(p\otimes p)\Delta(x^2)=2x\otimes x\neq0$ in the ordinary tensor square.
It is therefore not a coalgebra quotient there.  The surviving term has
total height two.  Discarding that total height in the tensor target
preserves the coproduct identity through height one.

\subsection{The total-height tensor}
\label{sec:height-tensor}

Let $\mathcal V$ be ordinary vector spaces, super vector spaces, or cochain
complexes over $\kk$.  In the super case morphisms preserve parity.
In the cochain case the differential has degree $+1$, and morphisms are
degree-zero chain maps.  A $Q_+$-graded object is a family
$(V_\beta)_{\beta\in Q_+}$ in $\mathcal V$, and its morphisms preserve charge.
For complexes this means $d(V_\beta^j)\subseteq V_\beta^{j+1}$.
For homogeneous vectors, write $p(a)\in\mathbb Z/2$ for super parity and
$|a|\in\mathbb Z$ for cohomological degree.  The symmetries are
\[
 \sigma(a\otimes b)=
 \begin{cases}
 b\otimes a,&\text{ordinary vector spaces},\\
 (-1)^{p(a)p(b)}b\otimes a,&\text{super vector spaces},\\
 (-1)^{|a||b|}b\otimes a,&\text{cochain complexes}.
 \end{cases}
\]
For complexes the tensor differential is
$d(a\otimes b)=da\otimes b+(-1)^{|a|}a\otimes db$.
Height is a function of charge, independent of these sign conventions.
In particular, an even vector can have height one.

For $N\geq0$, let $\mathsf{Gr}_{\leq N}(Q_+;\mathcal V)$ consist of such
objects supported in heights at most $N$.  Define
\begin{equation}
 (V\mathbin{\bar\otimes_N}W)_\beta
 =\bigoplus_{\gamma+\delta=\beta}V_\gamma\otimes W_\delta,
 \qquad \htop(\beta)\leq N.\label{finite-tensor}
\end{equation}
The tensor unit is $\kk$ in charge zero, with even parity or cohomological
degree zero and zero differential as appropriate.  Write
$\mathsf{Gr}_{\leq N}(Q_+)$ for the ordinary case.
The super and cochain constructions extend the tensor formalism.
The Borcherds recognition statements use the ordinary category.

\begin{proposition}[Monoidal height restriction]\label{monoidal-restriction}
Formula~\eqref{finite-tensor}, with the stated unit and symmetry, defines
a symmetric monoidal category.  Truncation
\[
 \tau_{\leq N}V=\bigoplus_{\htop(\beta)\leq N}V_\beta
\]
from nonnegatively $Q_+$-graded objects is strong symmetric monoidal.
Restriction from height $N$ to $M\leq N$ is also strong symmetric
monoidal, and the restrictions satisfy
\begin{equation}
 \tau_{\leq M}\tau_{\leq N}=\tau_{\leq M}.\label{transition}
\end{equation}
Their monoidal comparisons agree under successive restrictions.
\end{proposition}

\begin{proof}
In charge $\beta$, both parenthesizations of a triple tensor product
identify with
\[
 \bigoplus_{\gamma+\delta+\epsilon=\beta}
 U_\gamma\otimes V_\delta\otimes W_\epsilon,
 \qquad \htop(\beta)\leq N.
\]
Nonnegative height ensures that every intermediate sum has height at
most $\htop(\beta)$.  Thus both parenthesizations retain exactly these
summands.  Use the associator of $\mathcal V$ on each summand.
The pentagon on four factors is the base pentagon on each retained
quadruple.  The unit maps are the base unit maps on the summands of
charge zero in the unit factor.  Their triangle holds summandwise.
The stated symmetry exchanges the two charge indices, and its hexagons
and square hold by the same argument.  In complexes all these maps
commute with the charge-preserving differential.

If $\gamma+\delta=\beta$ and $\htop(\beta)\leq N$, both inputs have
height at most $N$.  Consequently the canonical comparison
\[
 \tau_{\leq N}V\bar\otimes_N\tau_{\leq N}W
 \longrightarrow\tau_{\leq N}(V\otimes W)
\]
is the identity on every retained summand.  The unit comparison is also
the identity.  These comparisons commute with the associators, unit maps
and symmetries just described.  For $M\leq N$, the same argument gives
\[
 \tau_{\leq M}(V\bar\otimes_N W)
 \cong\tau_{\leq M}V\bar\otimes_M\tau_{\leq M}W.
\]
Both composites of comparisons under three successive bounds are the
identity on each retained summand.  This proves the final assertion.
\end{proof}

These categories are truncated in height, without a dimension bound on
their objects.  They are generally not finite tensor categories in the
rigid sense.  Indeed, if $V$ is nonzero and supported in one positive
charge, every charge-preserving evaluation $V\bar\otimes_N W\to\kk$
is zero.  It cannot satisfy the duality triangle for $V$.

\subsection{Bialgebra identities in a window}

A bialgebra object $C$ has charge-preserving maps
\[
 m:C\bar\otimes_N C\to C,\quad \eta:\kk\to C,\quad
 \Delta:C\to C\bar\otimes_N C,\quad \varepsilon:C\to\kk.
\]
With associators and unit identifications understood, the algebra and
coalgebra axioms are
\begin{align*}
 m(m\bar\otimes_N1)&=m(1\bar\otimes_N m),&
 m(\eta\bar\otimes_N1)&=1_C=m(1\bar\otimes_N\eta),\\
 (\Delta\bar\otimes_N1)\Delta&=(1\bar\otimes_N\Delta)\Delta,&
 (\varepsilon\bar\otimes_N1)\Delta&=1_C
   =(1\bar\otimes_N\varepsilon)\Delta.
\end{align*}
The four compatibility identities are
\begin{align}
 \Delta m&=(m\bar\otimes_N m)
 (1\bar\otimes_N\sigma\bar\otimes_N1)
 (\Delta\bar\otimes_N\Delta),\label{bialgebra-identities}\\
 \Delta\eta&=\eta\bar\otimes_N\eta,\notag\\
 \varepsilon m&=\varepsilon\bar\otimes_N\varepsilon,
 \qquad \varepsilon\eta=1_\kk.\notag
\end{align}
For example, coassociativity on three specified output charges reads
\[
 (\Delta_{\alpha,\beta}\otimes1)\Delta_{\alpha+\beta,\gamma}
 =(1\otimes\Delta_{\beta,\gamma})\Delta_{\alpha,\beta+\gamma},
 \qquad \htop(\alpha+\beta+\gamma)\leq N.
\]
In complexes every structure map is a chain map.  Thus $d1=0$,
$\varepsilon d=0$, $d$ is a derivation, and
$\Delta d=d_{C\bar\otimes_N C}\Delta$ is the coderivation identity.

Strong monoidality transports all these identities under truncation.
It likewise transports Lie and associative algebra structures separately.
For a nonnegatively graded Lie algebra $\mathfrak g$,
$\tau_{\leq N}U(\mathfrak g)$ is therefore a bialgebra object in the
height category.  Its underlying ordinary vector space has two distinct
descriptions: it is the algebra quotient by heights above $N$, and it
embeds as a subcoalgebra of $U(\mathfrak g)$.
The quotient projection need not preserve the ordinary coproduct.
The algebra quotient and subcoalgebra structures need not form an
ordinary bialgebra together.

\begin{remark}[Two successive polynomial windows]\label{polynomial-window}
Put $C_N=\bigoplus_{n=0}^N\kk x^n$, with $|x|=0$,
$\deg_Qx=\alpha_i$, and $d=0$.  Polynomial multiplication where the total
height is at most $N$, together with
\[
 \Delta_N(x^n)=\sum_{r+s=n}\binom nr x^r\otimes x^s,
 \qquad \varepsilon_N(x^n)=\delta_{n0},\qquad \eta_N(1)=1,
\]
makes $C_N$ a bialgebra object by
Proposition~\ref{monoidal-restriction}.
At height one, $x\otimes x$ is absent from $C_1\bar\otimes_1C_1$.
At height two it occurs with coefficient two in
\eqref{polynomial-cross-term}.  The tensor restriction
\[
 Q:C_2\bar\otimes_2C_2\longrightarrow C_1\bar\otimes_1C_1
\]
discards all total-height-two terms.  For $p:C_2\to C_1$, this gives
$Q\Delta_2=\Delta_1p$ on $1,x,x^2$.
The ordinary rectangular projection $p\otimes p$ instead retains
$2x\otimes x$ on $x^2$, while $\Delta_1p(x^2)=0$.
\end{remark}

\begin{remark}[Unit normalization]\label{unit-normalization}
Multiplicativity of $\Delta$ does not imply the other three compatibility
identities.  Concentrate $C=\kk e\oplus\kk f$ in charge zero and set
\[
 e^2=e,\quad f^2=f,\quad ef=fe=0,\quad 1=e+f,
\]
\[
 \Delta e=e\otimes e,\quad \Delta f=f\otimes f,
 \qquad\varepsilon(e)=\varepsilon(f)=1.
\]
Products of basis idempotents vanish unless their indices agree, which
proves associativity and multiplicativity of $\Delta$.
The diagonal coproduct is coassociative, and both counit identities hold
on $e,f$.  Multiplication by $e+f$ gives both unit identities.  Nevertheless,
\[
 \Delta1=e\otimes e+f\otimes f\neq1\otimes1,\qquad
 \varepsilon(1)=2,\qquad
 \varepsilon(ef)=0\neq\varepsilon(e)\varepsilon(f).
\]

Suppose instead that $C_0=\kk1$ and all maps preserve the nonnegative
charge grading.  Assume associativity, coassociativity, both unit
identities, both counit identities and multiplicativity of $\Delta$.
Write $\Delta1=c(1\otimes1)$ and $\varepsilon(1)=a$.
The counit identity gives $ca=1$, while multiplicativity gives $c^2=c$.
Thus $c=a=1$.  The counit vanishes in positive charge, so
$\varepsilon m=\varepsilon\otimes\varepsilon$ holds there and on
$C_0\otimes C_0$.  This restores all omitted normalizations.
Here $C_0$ means charge zero, and in the cochain case it is the complex
$\kk$ concentrated in cohomological degree zero.
A chosen unit line, or a statement about $H^0(C)$, does not give this
hypothesis.
\end{remark}
